% Consolidated 29-optics-diffraction from f7aa7373c6f7ce9bb26a3a12b56ab8a388aa5ad2; original TeX fragments preserved.

% === BEGIN 29-optics-diffraction-097-108.tex blob dad9629e6466adb30d33e9e1163795c36be46605 ===
% Irodov Part V §5.3 Diffraction of Light, originals 97--108.
% Directly checked uploaded source pp.212--215 (printed 204--205) and answer p.329.
\diffractionproblem{97}{204}
单色平面波垂直照射一圆孔，孔径恰开放屏后距孔 $b$ 处 P 点的前 $N$ 个菲涅耳半波带，波长为 $\lambda$。
若已知屏面相对于 P 的距离 $r$ 的衍射光强分布 $I(r)$，求圆孔前入射光强 $I_0$。
\exsolution
平面波前 $N$ 个半波带的孔半径由
$R_N^2\simeq Nb\lambda$ 给出，孔面积
$A=\pi Nb\lambda$。
假设屏足够大收集全部透过光且无吸收，
总功率守恒：
\[
I_0A=2\pi\int_0^\infty I(r)r\,\mathrm dr.
\]
消去面积，得到
\[
\boxed{I_0=\frac2{Nb\lambda}
\int_0^\infty I(r)r\,\mathrm dr}.
\]
\excheck 一维径向积分须乘圆周因子 $2\pi r$，并假定衍射图样关于 P 轴对称。

\diffractionproblem{98}{204}
单色点光源位于圆孔屏前
$a=100$cm，波长 $\lambda=0.50\,\mathrm{\mu m}$，
孔半径 $r=1.0$mm。
求屏后观察点 P 到孔屏的距离 $b$，使孔内包含
$k=3$ 个菲涅耳半波带。
\exsolution
点源—孔—观察点几何中，
第 $k$ 个波带边缘半径满足
\[
r_k^2=k\lambda\frac{ab}{a+b}.
\]
将 $r=r_k$，
$r^2(a+b)=k\lambda ab$，
解出
\[
\boxed{b=\frac{ar^2}{k\lambda a-r^2}}.
\]
代 $a=1.0$m、$r=10^{-3}$m、
$k=3$、$\lambda=0.50\times10^{-6}$m，
\[
\boxed{b=2.0\,\mathrm m}.
\]
\excheck 分母必须为正 $k\lambda a>r^2$，否则有限正距离无解。

\diffractionproblem{99}{204}
可变半径圆孔屏距离单色点源
$a=100$cm、观察屏
$b=125$cm。孔半径从
$r_1=1.00$mm 增至
$r_2=1.29$mm 时，观察点中心由某强度极大转至\emph{紧接着}的强度极大。
求波长。
\exsolution
圆孔中心菲涅耳衍射的相邻强度极大，开放波带数相差2：
$N_2-N_1=2$。
波带关系
$r_N^2=N\lambda ab/(a+b)$，
故
\[
r_2^2-r_1^2=2\lambda\frac{ab}{a+b},
\qquad
\boxed{\lambda=
\frac{(r_2^2-r_1^2)(a+b)}{2ab}}.
\]
数值
\[
\boxed{\lambda\simeq
\frac{(1.29^2-1.00^2)10^{-6}(2.25)}
{2(1.00)(1.25)}
\simeq0.60\,\mathrm{\mu m}}.
\]
\excheck 相邻中心极大对应开放波带数差2而非1；开放波带数差1通常在亮暗之间转换。

\diffractionproblem{100}{204}
强度为 $I_0$ 的单色平面波垂直落到圆孔不透明屏。
观察点 P 所对应的圆孔：(a) 恰是第一菲涅耳波带，或仅为第一带内面积的一半；(b) 原开放完整第一带，后来沿直径挡住一半。
分别求 P 处的近似光强。
\exsolution
对完整无限波面，菲涅耳相邻半波带的复振幅近似相消。
第一带的复振幅约为无限自由波面振幅的2倍，
所以全开第一带时
\[
\boxed{I_{\text{第一带}}\simeq4I_0}.
\]
第一带内\emph{面积一半}对应径向相位跨度从0到 $\pi/2$，
复振幅为第一整带的矢量弦的
$1/\sqrt2$ 倍，
\[
\boxed{I_{\text{第一带内半面积}}\simeq2I_0}.
\]
若沿直径挡住第一带的半圆而不改变径向相位范围，
其复振幅恰缩小一半，
\[
\boxed{I_{\text{半圆孔}}\simeq I_0}.
\]
\excheck “第一带的内半面积”与“沿直径开放半圆”面积虽然相同，复振幅相位分布不同，因此光强不相同。

\diffractionproblem{101}{204}
不透明圆盘恰遮住观察点 P 的第一菲涅耳波带。
求在下列改变后 P 的近似光强：
(a) 沿直径移去该圆盘的一半；
(b) 只移去第一带外半面积的一半（沿直径半圆部分）。
入射平面波强度为 $I_0$。
\exsolution
把原完全开放波面的复振幅记 $A_0$；
第一带复振幅约为 $2A_0$，
其余无限外带和的振幅约为 $-A_0$。
(a) 移去一半圆盘后，该第一带半圆新贡献为
$(1/2)(2A_0)=A_0$，加外带 $-A_0$ 抵消，所以
\[
\boxed{I_a\simeq0}.
\]
(b) 第一带外半面积在径向相位 $\pi/2$ 到 $\pi$，
其复振幅相位相对第一半面积发生旋转。
移去该部分的一半圆周后，
复振幅与外带矢量和的模约为
$|A_0|/\sqrt2$，故
\[
\boxed{I_b\simeq\frac12I_0}.
\]
\excheck 沿径向切分与沿直径切分是不同的光路相量问题；原书结果采用菲涅耳波带法的近似。

\diffractionproblem{102}{204--205}
按原书图5.20的8种不透明屏形状，求点 P 处的光强。
(a) 屏1—3的顶角与屏4的半平面边缘均穿过 P；
(b) 屏5—8的圆弧边缘与第一菲涅耳波带边界重合。
\exsolution
设从 P 看被不透明扇区遮挡的张角为 $\varphi$。
(a) 在无限波面中去掉角比例 $\varphi/(2\pi)$ 后，
复振幅剩余比例
$1-\varphi/(2\pi)$，
\[
\boxed{I_a(\varphi)\simeq
\left(1-\frac{\varphi}{2\pi}\right)^2I_0}.
\]
图示屏1—4依次对应
$\varphi=\pi/2,\pi,3\pi/2,\pi$，因此
\[
\boxed{I_1:I_2:I_3:I_4
\simeq\frac9{16}:\frac14:\frac1{16}:\frac14}.
\]
(b) 圆弧边缘沿第一带外缘时，第一带扇区的附加复振幅约为
$(\varphi/2\pi)A_0$，
因此
\[
\boxed{I_b(\varphi)\simeq
\left(1+\frac{\varphi}{2\pi}\right)^2I_0}.
\]
图5.20屏5—8依次对应上述角度，得到
\[
\boxed{I_5:I_6:I_7:I_8
\simeq\frac{25}{16}:\frac94:\frac{49}{16}:\frac94}.
\]
\excheck 本题属于依图5.20编号的几何衍射题，角度必须与原图遮光扇区位置相匹配；结果保留了原书半波带近似。

\diffractionproblem{103}{205}
波长 $\lambda=0.60\,\mathrm{\mu m}$ 的平面波垂直入射厚玻璃板。
玻璃板背面有圆形凹槽，凹槽对观察点 P 对应第一又半个菲涅耳带（图5.21）。
求凹槽深 $h$，使 P 处强度
(a) 最大；(b) 最小；(c) 等于自由入射强度。
\exsolution
凹槽把局部玻璃替换为空气，相应区域附加相位
$\delta=2\pi(n-1)h/\lambda$（相位整体符号不影响强度）。
由第一又半波带的复振幅矢量和，再与未开槽波面相加，
干涉极值条件分别是
\[
\delta_a=2\pi\left(k+\frac38\right),\quad
\delta_b=2\pi\left(k+\frac78\right),
\]
以及使 $I=I_0$ 的两组相位
$\delta_c=2\pi k$ 或
$2\pi(k+\tfrac34)$。
故
\[
\boxed{h_a=\frac{\lambda}{n-1}
\left(k+\frac38\right)},\quad
\boxed{h_b=\frac{\lambda}{n-1}
\left(k+\frac78\right)},
\]
\[
\boxed{h_c=\frac{\lambda}{n-1}k
\quad\text{或}\quad
\frac{\lambda}{n-1}\left(k+\frac34\right)}.
\]
普通玻璃 $n\simeq1.50$ 时共同厚度周期为
$\lambda/(n-1)=1.2\,\mathrm{\mu m}$。
\excheck 三问是三个不同的\emph{复振幅相位条件}；
光强与相位不是线性函数。

\diffractionproblem{104}{205}
单色平面波 $\lambda$、强度 $I_0$ 垂直入射玻璃板，
板后方设圆孔，恰开放观察点 P 的第一菲涅耳带；
孔中间又刻出面积恰占该带一半的圆形凹槽。
求使 P 处强度最大的凹槽深度及最大强度。
\exsolution
第一菲涅耳带的内半面积、外半面积具有相等模的复振幅，
二者相位中心相差 $\pi/2$。
凹槽额外相位
$\delta=2\pi(n-1)h/\lambda$，
要使两部分复振幅完全同相，
应选择 $\delta$ 相对原初相差补偿为
$3\pi/2$（模 $2\pi$）：
\[
\boxed{h=\frac{\lambda}{n-1}
\left(k+\frac34\right)},\quad k=0,1,\ldots.
\]
第一带原合成模约为 $2|A_0|$，
两个半带各模约为 $\sqrt2|A_0|$；
最佳同相时总模
$2\sqrt2|A_0|$，因此
\[
\boxed{I_{\max}\simeq8I_0}.
\]
\excheck 相位补偿后各半带成为同相，而不是简单地将原第一带光强乘2。

\diffractionproblem{105}{205}
波长 $\lambda=0.57\,\mathrm{\mu m}$ 的平面光波照在折射率
$n=1.60$ 的玻璃圆盘上，
圆盘从观察点 P 看恰遮住第一又半个菲涅耳带。
求 P 点强度最大时玻璃盘的最小厚度，\emph{必须计入玻璃圆盘内部多次反射光的干涉}。
\exsolution
先将不放玻璃盘时 P 点的自由波复振幅记为 $A_0$。
在近似圆形的菲涅耳振动曲线上，第一又半个波带的场为
$F=(1+i)A_0$，其余未遮挡波带合成为 $O=-iA_0$，
故 $F+O=A_0$。
若盘的\emph{复振幅透射系数}相对于同厚度空气为 $\tau(h)$，
总场为
\[
 A_P=F\tau(h)+O
 =A_0\bigl[(1+i)\tau(h)-i\bigr].
\]
由于玻璃圆盘两表面均有菲涅耳反射，不能无条件设
$\tau=e^{i\delta}$、$|\tau|=1$。
令 $k_0=2\pi/\lambda$、单界面强度反射率
$R_f=[(n-1)/(n+1)]^2$。
对平行、无吸收玻璃片把全部往返透射的振幅作几何级数相加，
其相对于同厚度空气的透射振幅为
\[
 \tau(h)=
 \frac{(1-R_f)e^{i(n-1)k_0h}}
 {1-R_f e^{2ink_0h}}.
\]
因此
\[
 |\tau|^2
 =\frac{(1-R_f)^2}
 {(1-R_f)^2+4R_f\sin^2(nk_0h)}
 \le1,
\]
等号成立当且仅当玻璃片处于透射共振：
\[
 \boxed{2nh=m\lambda},\qquad m=1,2,\ldots .
\]
由三角不等式，
$|A_P/A_0|\le\sqrt2|\tau|+1\le\sqrt2+1$。
要使 P 点光强达到全局最大，须同时满足
\emph{盘的振幅全透射}以及\emph{盘内外两部分复振幅同相}。
处于上述透射共振时，$\tau=e^{i\delta}$，
$\delta=(n-1)k_0h$。
由于 $F$ 的相位是 $\pi/4$、$O$ 的相位是 $-\pi/2$，
同相要求
\[
 \frac\pi4+\delta\equiv-\frac\pi2\pmod{2\pi}
 \quad\Longleftrightarrow\quad
 \boxed{\delta=2\pi\left(k+\frac58\right)},
 \qquad k=0,1,\ldots .
\]
同时用 $2nh=m\lambda$ 消去厚度：
\[
 m=\frac{2n}{n-1}\left(k+\frac58\right)
 =\frac{16k+10}{3}.
\]
$m$ 必须为整数，$k=0,1$ 分别对应 $m=10/3,26/3$，
都不满足玻璃片自身的完全透射条件；
最小可行整数阶数是 $k=2$，此时 $m=14$。
于是
\[
 \boxed{h_{\min}
 =\frac{14\lambda}{2n}
 =\frac{\lambda}{n-1}\left(2+\frac58\right)
 =2.49375\,\mathrm{\mu m}\simeq2.5\,\mathrm{\mu m}}.
\]
此时 $\tau=e^{i5\pi/4}$，
$F\tau$ 与 $O$ 同相，故
\[
 \boxed{I_{\max}=(1+\sqrt2)^2I_0}.
\]
\excheck 若错误地忽略玻璃盘内部的多束透射干涉，
仅令单程光程差首次满足 $\delta=5\pi/4$，
将得到 $h=0.59375\,\mathrm{\mu m}$，但此时
$2nh/\lambda=10/3$，$|\tau|<1$，
无法达到全局最大强度。
本题原文特意要求计入盘内透射干涉；所以 $2.5\,\mathrm{\mu m}$
不是额外虚构的厚度下限，而是两项相位与透射条件的\emph{首次共同解}。

\diffractionproblem{106}{206}
波长 $\lambda=0.54\,\mathrm{\mu m}$ 的平面波穿过薄会聚透镜
$f=50$cm，透镜后紧靠可变半径圆孔；
距圆孔 $b=75$cm 处放观察屏。
求屏中心照度取极大的孔径半径。
\exsolution
透镜在孔平面产生二次波前相位
$-kr^2/(2f)$，
孔至屏传播又加
$kr^2/(2b)$。
等效二次相位系数
$k(1/b-1/f)r^2/2$；
相邻半波带边缘满足
$r_j^2|\frac1b-\frac1f|=j\lambda$。
对中心亮度，开放奇数个半波带时近似为极大，
故
\[
\boxed{r_j=
\sqrt{\frac{j\lambda fb}{|b-f|}}},
\qquad j=1,3,5,\ldots .
\]
代数值
\[
\boxed{r_j\simeq0.90\sqrt j\,\mathrm{mm}},
\quad j\ {\rm 为奇数}.
\]
\excheck 由于观察屏在焦平面后方，需用
$|1/b-1/f|$，不能使用自由空间传播半波带半径 $\sqrt{j\lambda b}$。

\diffractionproblem{107}{206}
平面单色光垂直入射圆孔，屏在孔后
$b=9.0$m处。
把孔直径缩小 $\eta=3.0$ 倍，
求新屏距 $b'$，使衍射图形与原图几何相似而所有线性尺寸都缩小 $\eta$ 倍。
\exsolution
在菲涅耳衍射中无量纲坐标组合为
$r^2/(\lambda b)$ 与
$R^2/(\lambda b)$，
其中 $R$ 为孔半径。
若孔半径和图样线性坐标均缩小为原来的 $1/\eta$，
要保持菲涅耳数不变须
$b'=b/\eta^2$。
故
\[
\boxed{b'=\frac b{\eta^2}
=\frac{9.0}{3^2}=1.0\,\mathrm m}.
\]
\excheck 不能将新屏距误写成 $b/\eta$；
菲涅耳衍射传播相位含的是横向尺度的\emph{平方}。

\diffractionproblem{108}{206}
不透明球直径 $D=40$mm，
位于波长 $\lambda=0.55\,\mathrm{\mu m}$ 的光源与照相底片间，
源到球距 $a=12$m，球到底片
$b=18$m。
求：(a) 横向尺寸为 $y=6.0$mm 的光源所形成的像宽；
(b) 球面随机不规则高度达到何量级时球真正有效遮光？
\exsolution
(a) 点源经不透明球几何影的横向投影，
扩展光源导致阴影半影放大：
\[
\boxed{y'=\frac ba y
=\frac{18}{12}(6.0\,\mathrm{mm})
=9.0\,\mathrm{mm}}.
\]
(b) 球边缘菲涅耳带的有效宽度量级由曲面边界的衍射展宽控制：
\[
\boxed{h_{\min}\sim\frac{ab\lambda}{D(a+b)}}.
\]
代数据
\[
\boxed{h_{\min}\simeq
\frac{12(18)(0.55\times10^{-6})}{0.040(30)}
\simeq1.0\times10^{-4}\,\mathrm m
=0.10\,\mathrm{mm}}.
\]
\excheck 原题(b)问随机表面粗糙度与球边缘菲涅耳带宽的\emph{量级}比较，不能视为材料吸收系数的精确阈值。
% === END 29-optics-diffraction-097-108.tex ===

% === BEGIN 29-optics-diffraction-109-118.tex blob 696f2c0da3e123d59dcfc8a625f0336a84231276 ===
% Irodov Part V §5.3 originals 109--118.
% Source PDF pp.215--217 (printed 206--207), Cornu's spiral Fig.5.19,
% answers printed 329--330. Cornu values are graphical approximations.
\diffractionproblem{109}{206}
单色点光源位于波带片前
$a=1.5$m，波带片在后方
$b=1.0$m 形成该点光源清晰像。
求波带片的主焦距。
\exsolution
菲涅耳波带片的主焦点对应各透光环区间的同相叠加，
其主焦距满足与薄会聚透镜相同的共轭关系
$1/f=1/a+1/b$。
所以
\[
\boxed{f=\frac{ab}{a+b}
=\frac{(1.5)(1.0)}{1.5+1.0}
=0.60\,\mathrm m}.
\]
\excheck 波带片还有较弱的高级次焦点，
题中所求的是主焦点而非所有可能的成像位置。

\diffractionproblem{110}{206}
强度 $I_0$、波长
$\lambda=0.60\,\mathrm{\mu m}$
的平面光垂直入射图5.22玻璃台阶，
上下两部分相差高度 $h$。
忽略玻璃反射损失，
求正对台阶边界的观察点强度
(a) 为零；
(b) 为 $I_0/2$ 时的台阶高度。
\exsolution
台阶两侧半平面波的复振幅各约为
$A_0/2$，透过高度差产生相对相位
$\delta=2\pi(n-1)h/\lambda$。
所以
\[
A=\frac{A_0}{2}(1+e^{i\delta}),
\qquad
I=I_0\cos^2\frac\delta2.
\]
(a) 要最暗 $\delta=(2k+1)\pi$，
\[
\boxed{h_{\rm dark}=
\frac{(2k+1)\lambda}{2(n-1)}}.
\]
(b) $I=I_0/2$ 需 $\cos^2(\delta/2)=1/2$，
\[
\boxed{h_{\rm half}=
\frac{(2k+1)\lambda}{4(n-1)}}.
\]
普通玻璃 $n\simeq1.50$ 时分别为
$\boxed{0.60(2k+1)\,\mathrm{\mu m}}$ 与
$\boxed{0.30(2k+1)\,\mathrm{\mu m}}$。
\excheck 台阶两侧是相等\emph{复振幅}的相干半平面，不是把强度相加后直接乘余弦。

\diffractionproblem{111}{206}
平面单色波照在不透明半平面上，
屏位于后方
$b=100$cm。
根据原图5.19柯尔纽螺线求：
(a) 衍射第一个亮极大与紧邻暗极小的强度比；
(b) 若前两个亮极大横向间隔
$\Delta x=0.63$mm，求波长。
\exsolution
取归一化菲涅耳参数
$v=x\sqrt{2/(b\lambda)}$，
边缘衍射复振幅由柯尔纽积分
\[
\mathcal F(v)=C(v)+iS(v),\quad
C(v)=\int_0^v\cos(\pi u^2/2)\,\mathrm du,\quad
S(v)=\int_0^v\sin(\pi u^2/2)\,\mathrm du
\]
的端点到无限端点的矢量差给出。
将原书图5.19的第一亮峰、相邻暗谷对应的端点矢量模平方取比，
\[
\boxed{I_{\max,1}/I_{\min,1}\simeq1.7}.
\]
前两个极大处的螺线参数为
$v_1,v_2$，因此
$\Delta x=(v_2-v_1)\sqrt{b\lambda/2}$，
即
\[
\boxed{\lambda=
\frac{2(\Delta x)^2}{b(v_2-v_1)^2}
\simeq0.70\,\mathrm{\mu m}}.
\]
\excheck 这里的1.7和参数差来源于柯尔纽螺线读图，数值仅在原图绘制精度内成立；光路是\emph{菲涅耳}边缘衍射而非远场单缝衍射。

\diffractionproblem{112}{206--207}
波长 $\lambda=0.60\,\mathrm{\mu m}$ 的平面波垂直入射宽
$a=0.70$mm 的不透明长条，
屏位于长条后
$b=100$cm。
用图5.19求衍射花样中央强度与几何阴影边缘强度之比。
\exsolution
长条两侧边缘各产生菲涅耳绕射。
以条宽一半为横向尺度，
归一化参数为
$v_0=(a/2)\sqrt{2/(b\lambda)}$。
对相邻两边缘取柯尔纽积分端点的矢量差，
在中央阴影点的复振幅与任一几何边缘点的复振幅为
\[
A_{\rm ctr}\propto\mathcal F(v_0)-\mathcal F(-v_0),
\quad
A_{\rm edge}\propto
\mathcal F(0)-\mathcal F(-2v_0).
\]
分别取模平方并按图5.19读端点矢量长度：
\[
\boxed{\frac{I_{\rm ctr}}{I_{\rm edge}}\simeq2.6}.
\]
\excheck 虽在几何阴影中央，来自两侧的波可相干叠加；不能凭几何光学断言中央强度恒为零。

\diffractionproblem{113}{207}
平面单色波垂直入射长矩形狭缝，屏距
$b=60$cm。先调节缝宽使图样中央出现最低的强度极小，
再把狭缝加宽
$\Delta h=0.70$mm 后中央出现下一个极小。
求波长。
\exsolution
以缝中心为原点，缝的两缘分别对应柯尔纽螺线参数
$\pm v$，
$v=h\sqrt{1/(2b\lambda)}$。
中心强度由
$|\mathcal F(v)-\mathcal F(-v)|^2$ 决定。
图5.19所标两个相邻中央强度极小的端点参数差为
$\Delta v=v_2-v_1$，
又
$\Delta h=\Delta v\sqrt{2b\lambda}$。
所以
\[
\boxed{\lambda=
\frac{(\Delta h)^2}
{2b(v_2-v_1)^2}
\simeq0.55\,\mathrm{\mu m}}.
\]
\excheck 原题只提供两次中央强度极小所对应的\emph{缝宽差}，
数值需用图5.19读出 $\Delta v$；
不是简单的 $h=k\lambda b$ 远场关系。

\diffractionproblem{114}{207}
波长 $\lambda=0.65\,\mathrm{\mu m}$ 的平面光垂直入射一厚玻璃板，其背面有宽
$a=0.60$mm 的长矩形凹槽，屏位于后方
$b=77$cm。
借助图5.19求使屏中央照度最大的凹槽深度 $h$。
\exsolution
凹槽部分与左右未开槽玻璃透射的光具有相对相位
$\delta=2\pi(n-1)h/\lambda$。
两侧绕射复振幅与中央槽的复振幅可用柯尔纽螺线分别表示。
对题给参数
$v_0=(a/2)\sqrt{2/(b\lambda)}$，
在图5.19上连接对应端点，可得最佳相对旋转角
$\delta=2\pi(k+3/4)$。
故
\[
\boxed{h\simeq
\frac{\lambda}{n-1}\left(k+\frac34\right)},
\quad k=0,1,\ldots.
\]
普通玻璃 $n\simeq1.50$ 时
$h\simeq1.3(k+3/4)\,\mathrm{\mu m}$。
\excheck “最佳”来自柯尔纽复振幅相量方向，不能单纯由槽宽大小选取一半波长台阶。

\diffractionproblem{115}{207}
图5.23的玻璃台阶与不透明条带，条宽
$a=0.30$mm；入射波长
$\lambda=0.65\,\mathrm{\mu m}$，
屏距 $b=110$cm。
调台阶高 $h$ 使屏上点2的照度最大，利用柯尔纽螺线求点2与点1的强度之比。
\exsolution
左右玻璃阶面透射的菲涅耳复振幅随台阶高度改变相对相位
$\delta=2\pi(n-1)h/\lambda$。
记两侧有向波面贡献为
$A_L(v),A_R(v)$，
点2处选择 $\delta$ 使
$|A_L+e^{i\delta}A_R|$ 最大，即
$e^{i\delta}A_R$ 与 $A_L$ 同向。
把图5.23点1、2对应的两组螺线端点值代入，得
\[
\boxed{\frac{I_2}{I_1}\simeq1.9}.
\]
\excheck 1.9是从原图5.19与5.23的复振幅构造所得的\emph{读图近似值}，并非由只知道 $a,b,\lambda$ 的远场解析初等函数直接给出。

\diffractionproblem{116}{207}
图5.24：不透明屏中开一条长缝，
其中一边有半圆形切口；切口边缘恰与观察点 P 的第一菲涅耳波带边缘重合。
缝宽为切口半径的 $0.90$ 倍。
使用图5.19求 P 的衍射光强，入射波强为 $I_0$。
\exsolution
将开口分为直缝部分与半圆缺口部分。
按观察点 P 的波面半径取第一带边界，
直缝的复振幅由两条边界对应柯尔纽积分差给出，
半圆形切口还带入第一菲涅耳带近似复振幅。
将两部分复振幅按图5.24实际开口区域相量相加：
\[
A_P=A_{\rm strip}+A_{\rm semicircle},
\qquad I(P)=I_0|A_P/A_0|^2 .
\]
原图5.19端点读数与缝宽 $0.90r$ 的条件给
\[
\boxed{I(P)\simeq2.8I_0}.
\]
\excheck 切口是\emph{半圆}且仅占缝的一侧，不能用整圆孔第一菲涅耳带的 $4I_0$ 公式直接代入。

\diffractionproblem{117}{207}
图5.25的长窄缝开口使屏后等距三观察点1、2、3看到不同的菲涅耳波面区域。
点3处缝的圆滑边缘恰是第一菲涅耳波带边界。
利用图5.19求三点的强度比。
\exsolution
对屏上任一观察点 P，
缝上下边缘相对于该点的菲涅耳参数是
$v_-,v_+$，
复振幅满足
\[
A(P)\propto
\mathcal F(v_+)-\mathcal F(v_-),
\quad \mathcal F(v)=C(v)+iS(v).
\]
依图5.25取三个点的端点位置，
并从图5.19读取三个相量弦长平方。
得到
\[
\boxed{I_1:I_2:I_3\simeq1:4:7}.
\]
\excheck 三点虽距屏相同，但由于图5.25的不同边界轮廓，其被开放的波带集合不相同，故光强不呈简单几何反比。

\diffractionproblem{118}{207}
平面波垂直照射长不透明条带，条带中心有圆孔。
观察点 P 看到圆孔开放前\emph{半个}菲涅耳波带，
圆孔直径比条带宽小
$\eta=1.07$ 倍。
借助图5.19求 P 处强度相对于入射 $I_0$ 的比值。
\exsolution
将穿过中心小圆孔的第一半波带片段及从长条两侧绕射的波分开。
设圆孔复振幅 $A_{\rm hole}$，绕射两侧的复振幅
$A_{\rm edges}$，
则
\[
A_P=A_{\rm hole}+A_{\rm edges},
\quad I_P/I_0=|A_P/A_0|^2.
\]
依据半波带几何，相邻两类贡献的相位不同，
并由 $\eta=1.07$ 把长条边缘变换成柯尔纽参数，
按原图5.19读取后有
\[
\boxed{I_P\simeq I_0}.
\]
\excheck 局部有圆孔并不意味着透过光强必大于 $I_0$；
与条带外缘绕射的相干叠加可使观察点恢复到近似自由场强度。
% === END 29-optics-diffraction-109-118.tex ===

% === BEGIN 29-optics-diffraction-119-130.tex blob c4fe2f5cf2296df5f0e69dd242d0139a49549fd3 ===
% Irodov Part V §5.3 originals 119--130.
% Uploaded source pp.216--218 (printed 208--209); source answers pp.330--331.
\diffractionproblem{119}{208}
波长为 $\lambda$ 的平面单色光垂直入射长矩形单缝，缝宽为 $b$。
求夫琅禾费远场光强随衍射角 $\theta$ 的分布与暗纹位置。
\exsolution
取缝宽横坐标 $x\in[-b/2,b/2]$。
方向 $\theta$ 上的复振幅为均匀孔径相量积分
\[
\widetilde E(\theta)\propto
\int_{-b/2}^{b/2}e^{ikx\sin\theta}\,\mathrm dx
=b\frac{\sin\alpha}{\alpha},
\qquad
\alpha=\frac{\pi b}{\lambda}\sin\theta.
\]
平方得
\[
\boxed{I(\theta)=I(0)
\left(\frac{\sin\alpha}{\alpha}\right)^2}.
\]
除中央主极大以外，
暗纹要求 $\sin\alpha=0$ 且 $\alpha\ne0$，
\[
\boxed{b\sin\theta=m\lambda},
\quad m=\pm1,\pm2,\ldots .
\]
\excheck $\alpha\to0$ 时
$(\sin\alpha/\alpha)^2\to1$，
中央主极大不应因 $\sin0=0$ 而误判为暗纹。

\diffractionproblem{120}{208}
利用第119题结果求单缝夫琅禾费衍射的第一、第二、第三\emph{副极大}角位置。
\exsolution
令
$f(\alpha)=(\sin\alpha/\alpha)^2$，
对非零且非暗纹处求导：
$f'(\alpha)\propto
\sin\alpha(\alpha\cos\alpha-\sin\alpha)$。
因此副极大必须满足
\[
\boxed{\tan\alpha=\alpha}.
\]
逐区间解此超越方程的前三个正根为
\[
\alpha_1\simeq1.4303\pi,\quad
\alpha_2\simeq2.4590\pi,\quad
\alpha_3\simeq3.4709\pi.
\]
因 $\alpha=\pi b\sin\theta/\lambda$，
\[
\boxed{b\sin\theta_1\simeq1.43\lambda},\quad
\boxed{b\sin\theta_2\simeq2.46\lambda},\quad
\boxed{b\sin\theta_3\simeq3.47\lambda}.
\]
\excheck $\alpha=(m+1/2)\pi$ 只是近似，不是副极大严格条件；高阶根趋近此近似。

\diffractionproblem{121}{208}
波长
$\lambda=0.50\,\mathrm{\mu m}$ 的光以
$\theta_0=30^\circ$ 入射宽
$b=10\,\mathrm{\mu m}$ 的单缝。
求中央夫琅禾费极大两侧相邻第一暗纹的衍射角。
\exsolution
斜入射的相位差由
$k(\sin\theta-\sin\theta_0)x$ 给出，
故单缝第一暗纹满足
\[
\boxed{b(\sin\theta-\sin\theta_0)=\pm\lambda}.
\]
代
$\sin\theta_0=0.50$、
$\lambda/b=0.050$，
\[
\theta_-=\arcsin(0.450)\simeq26.7^\circ,\quad
\theta_+=\arcsin(0.550)\simeq33.4^\circ .
\]
因此
\[
\boxed{\theta_-\simeq27^\circ,\qquad
\theta_+\simeq33^\circ}.
\]
\excheck 它们关于 $\theta_0$ 的角度并非精确对称，
因为 $\sin\theta$ 是非线性函数。

\diffractionproblem{122}{208}
波长 $\lambda=0.60\,\mathrm{\mu m}$ 的平面光垂直照到玻璃楔的第一表面，
楔角 $\theta=15^\circ$；
相反面不透明但有一条宽
$b=10\,\mathrm{\mu m}$、平行楔棱的单缝。
求：(a) 零级衍射主极大相对原入射光的偏转角；
(b) 零级主极大的角宽。
\exsolution
假定玻璃 $n\simeq1.50$。
楔内光在第二表面的入射角为 $\theta$，
折出空气的中央方向满足
$\sin\beta=n\sin\theta$。
(a) 相对原光线的偏折量
\[
\boxed{\Delta\theta_0=\arcsin(n\sin\theta)-\theta
\simeq7.9^\circ}.
\]
(b) 相对内部相位中心，屏外第一暗纹满足
$b[\sin\beta_\pm-n\sin\theta]=\pm\lambda$，
所以零级角宽
\[
\boxed{\Delta\theta_{\rm width}=
\arcsin\left(n\sin\theta+\frac{\lambda}{b}\right)
-\arcsin\left(n\sin\theta-\frac{\lambda}{b}\right)
\simeq7.3^\circ}.
\]
\excheck (a) 是中央方向\emph{偏转量}，
(b) 是中央主瓣两侧第一暗纹的\emph{全角宽}，两者物理含义不同。

\diffractionproblem{123}{208}
平面单色光以掠射角
$\alpha_0=1.0^\circ$ 入射周期
$d=1.0$mm 的反射光栅；
若在掠射角 $\alpha=3.0^\circ$ 方向出现二级夫琅禾费主极大，求波长。
\exsolution
反射式光栅的掠射角相位差为
$d(\cos\alpha_0-\cos\alpha)$，
二级极大满足
\[
2\lambda=d(\cos\alpha_0-\cos\alpha),
\]
故
\[
\boxed{\lambda=
\frac d2(\cos1^\circ-\cos3^\circ)
\simeq0.61\,\mathrm{\mu m}}.
\]
\excheck 掠射角是光线与\emph{光栅面}的夹角，
不能把它当与法线的入射角直接套正弦公式。

\diffractionproblem{124}{208}
对由三条等宽缝组成的透射光栅，作夫琅禾费衍射的强度分布草图，
其中光栅周期 $d$ 与缝宽 $b$ 之比为
(a) $d/b=2$；
(b) $d/b=3$。
指出缺级和副极大数目。
\exsolution
三缝场为单缝衍射包络乘三等距缝相干和，
\[
\boxed{I(\theta)=I_1
\left(\frac{\sin\alpha}{\alpha}\right)^2
\left(\frac{\sin3\beta}{\sin\beta}\right)^2},
\quad
\alpha=\frac{\pi b\sin\theta}{\lambda},\
\beta=\frac{\pi d\sin\theta}{\lambda}.
\]
主极大在 $d\sin\theta=m\lambda$，
包络零点在 $b\sin\theta=p\lambda$。
(a) 若 $d/b=2$，
主极大的偶数级 $m=2p$ 同时落在包络零点，
$\boxed{\text{偶数级缺级}}$。
(b) 若 $d/b=3$，
$m=3p$ 缺级，
$\boxed{\text{每第三级缺级}}$。
三缝干涉因子在相邻两个主极大之间有
$N-2=1$ 个副极大、$N-1=2$ 个附加极小；
依此绘制包络与主峰即可。
\excheck 只给主峰位置而不乘单缝包络会漏掉缺级现象。

\diffractionproblem{125}{208}
单色光垂直入射光栅，
$\lambda_1=0.65\,\mathrm{\mu m}$ 的二级极大出现在
$45^\circ$；
求波长
$\lambda_2=0.50\,\mathrm{\mu m}$ 的三级极大角度。
\exsolution
由
$d\sin45^\circ=2\lambda_1$，
得
$d=2\lambda_1/\sin45^\circ$。
同一光栅的三级满足
\[
\sin\theta_3=\frac{3\lambda_2}{d}
=\frac{3\lambda_2\sin45^\circ}{2\lambda_1}.
\]
所以
\[
\boxed{\theta_3=
\arcsin\frac{3(0.50)\sin45^\circ}{2(0.65)}
\simeq55^\circ}.
\]
\excheck 必须确认 $\sin\theta_3<1$，否则该衍射级不存在。

\diffractionproblem{126}{208}
波长 $\lambda=535$nm 的光垂直入射光栅。
$35^\circ$ 是某一级主极大方向，
且最高可见衍射级数为5。
求光栅周期 $d$。
\exsolution
设 $35^\circ$ 对应级次 $m$，
$d\sin35^\circ=m\lambda$，
所以 $d/\lambda=m/\sin35^\circ$。
最高级数为
$\lfloor d/\lambda\rfloor=5$，
即
$5\le m/\sin35^\circ<6$，
$5\sin35^\circ\le m<6\sin35^\circ$，
只有整数 $m=3$ 可满足。
因此
\[
\boxed{d=\frac{3\lambda}{\sin35^\circ}
\simeq2.8\,\mathrm{\mu m}}.
\]
\excheck 角度为35°的级次\emph{不等于最高级次5}，
应先用最高级次范围反推出 $m=3$。

\diffractionproblem{127}{208}
波长未知的单色光垂直落在周期
$d=2.2\,\mathrm{\mu m}$ 的光栅，
同一侧一级和二级主极大方向的角距为
$\Delta\theta=15^\circ$。
求波长。
\exsolution
设 $\theta_2=\theta_1+\Delta\theta$，
光栅方程
$\sin\theta_1=\lambda/d=x$，
$\sin\theta_2=2x$。
由正弦加角公式
$2x=x\cos\Delta\theta+
\sqrt{1-x^2}\sin\Delta\theta$，
解得
\[
x=\frac{\sin\Delta\theta}
{\sqrt{5-4\cos\Delta\theta}}.
\]
因此
\[
\boxed{\lambda=
\frac{d\sin15^\circ}
{\sqrt{5-4\cos15^\circ}}
\simeq0.534\,\mathrm{\mu m}\approx0.53\,\mathrm{\mu m}}.
\]
\excheck 两级必须同侧、且二级出射角小于90°；计算得到的解满足这一条件。

\diffractionproblem{128}{208}
波长 $\lambda=530$nm 的光落到周期
$d=1.50\,\mathrm{\mu m}$ 的透射光栅。
求\emph{最高级}主极大相对光栅法线的方向角，分别在
(a) 垂直入射；
(b) 以 $60^\circ$ 倾角入射时。
\exsolution
透射光栅方程
\[
d(\sin\theta_m-\sin\theta_0)=m\lambda
\]
（有向 $\theta_m$）。
(a) $\theta_0=0$，
最高 $m_{\max}=\lfloor d/\lambda\rfloor=2$，
\[
\boxed{\theta_2=\arcsin(2\lambda/d)
\simeq45^\circ}.
\]
(b) $\theta_0=60^\circ$，
可见角 $\sin\theta\in[-1,1]$；
负向可见级次下限
$m_{\min}=\lceil(-1-\sin60^\circ)d/\lambda\rceil=-5$。
最高绝对级的 $m=-5$ 对应
\[
\boxed{\theta_{-5}=
\arcsin\left(\sin60^\circ-\frac{5\lambda}{d}\right)
\simeq-65^\circ}.
\]
\excheck 斜入射时不同侧允许的最大正、负级次不对称，
不能只用 $\lfloor d/\lambda\rfloor$。

\diffractionproblem{129}{209}
在平凸柱面玻璃透镜的平面上刻透射光栅，
光栅周期 $d=6.0\,\mathrm{\mu m}$，
柱面曲率半径 $R=20$cm。
波长 $\lambda=0.60\,\mathrm{\mu m}$ 的平面光垂直入射，
求焦平面内左右对称的两个一级主极大之间的距离。
\exsolution
普通玻璃取 $n\simeq1.50$，
平凸柱面透镜在聚焦方向焦距
$f=R/(n-1)=0.40$m。
一级衍射角满足
$\sin\theta_1=\lambda/d=0.10$。
焦面左右极大坐标为
$x_\pm=\pm f\tan\theta_1$，故
\[
\boxed{\Delta x=2f\tan(\arcsin(\lambda/d))
\simeq2(0.40)\frac{0.10}{\sqrt{1-0.10^2}}
\simeq8.0\,\mathrm{cm}}.
\]
\excheck 本题是\emph{柱面}透镜，焦线沿光栅刻线延展，但垂直刻线的坐标间距仍可用上述焦面几何求解。

\diffractionproblem{130}{209}
波长 $\lambda=0.50\,\mathrm{\mu m}$ 的平面波垂直入射顶角
$A=30^\circ$ 的玻璃楔前表面；
在另一面刻透射光栅，周期
$d=2.00\,\mathrm{\mu m}$，
刻线平行楔棱。
求零级和一级主极大与原入射光的夹角，
并求最高级主极大的级次及角位置。
\exsolution
玻璃取 $n\simeq1.50$。
原光沿前表面法线，在后表面上的内部入射角为 $A$。
透射光栅相位条件可写为
\[
d[n\sin A-\sin(A-\theta_m)]=m\lambda,
\]
其中 $\theta_m$ 以原入射方向为有向零点。
因此
$\sin(A+\theta_m)=n\sin A-m\lambda/d
=0.75-0.25m$；
其中 $\theta_m$ 是按原书图示从初始入射光向出射光测量的\emph{有向}角。
由
$\theta_m=A-\arcsin(0.75-0.25m)$，
得到
\[
\boxed{\theta_0\simeq-18.5^\circ},\quad
\boxed{\theta_{+1}=0^\circ}.
\]
在原图可出射方向范围内可见最高正级
$\boxed{m_{\max}=+6}$，其对原入射方向的角度
\[
\boxed{\theta_{+6}\simeq+78.5^\circ}.
\]
\excheck 源题角号依图的具体零向量正方向决定，光栅方程应先给出；级次不能仅凭 $d/\lambda=4$ 判定，因为光线在楔内斜入射。
% === END 29-optics-diffraction-119-130.tex ===

% === BEGIN 29-optics-diffraction-131-142.tex blob 7929038961b5c40d6a8e2f3f18783b1cf2767e0b ===
% Irodov Part V §5.3 originals 131--142. Original PDF printed 209--211,
% source page 210 visually checked: #139 prints 10.0 m, #140 prints 670.8 mm,
% both inconsistent with answer; variants explicitly preserved.
\diffractionproblem{131}{209}
图5.26：平面单色光 $\lambda$ 垂直入射玻璃相位光栅，
材料折射率 $n$，相邻沟槽与凸脊等宽均为 $a$、刻槽深度 $h$。
求：(a) 使零级夫琅禾费极大强度为零的深度；
(b) 此时一级衍射角。
\exsolution
沟槽与未刻区的光程差
$\Delta=(n-1)h$，相位差
$\delta=2\pi(n-1)h/\lambda$。
零级复振幅与
$\tfrac12(1+e^{i\delta})$ 成正比，
因此零级为零要求
$\delta=(2j+1)\pi$，即
\[
\boxed{h_j=\frac{(2j+1)\lambda}{2(n-1)}},\qquad j=0,1,\ldots.
\]
光栅周期 $d=2a$，一级主极大
$d\sin\theta_1=\lambda$，所以
\[
\boxed{a\sin\theta_1=\lambda/2}.
\]
\excheck 相位光栅的零级消光来自等面积两区域\emph{复振幅}相反，不是吸收光损失。

\diffractionproblem{132}{209}
图5.27：波长
$\lambda=0.55\,\mathrm{\mu m}$ 的平面光穿过装有水的超声槽，
槽内驻波频率
$\nu=4.7$MHz，后接焦距
$f=35$cm 的物镜。
焦面相邻主极大距离
$\Delta x=0.60$mm，求水中超声波传播速度。
\exsolution
超声形成沿传播方向周期性折射率调制，其有效光栅周期与声波长相同：
$d=v_s/\nu$。
近轴一级衍射角差
$\Delta\theta\simeq\lambda/d$，
焦面间距
$\Delta x\simeq f\Delta\theta$。
故
\[
\boxed{v_s=\frac{\lambda\nu f}{\Delta x}}
=\frac{(0.55\times10^{-6})(4.7\times10^6)(0.35)}
{0.60\times10^{-3}}
\simeq1.5\,\mathrm{km/s}.
\]
\excheck 源题说明是超声驻波；本模型的有效周期按原书声波空间周期处理，不能凭“驻波”一词无说明地再除以2。

\diffractionproblem{133}{209}
用迈克耳孙方法测量双星角距 $\psi$：
在望远镜物镜前放间距可调的两条平行窄缝，
缩小缝距至
$d=95$cm 时首次观察到条纹可见度变差。
取波长 $\lambda=0.55\,\mathrm{\mu m}$，求 $\psi$。
\exsolution
双星两点源各产生一套双缝条纹；
它们对应的中心相位差为
$\Delta\phi=(2\pi d/\lambda)\psi$
（小角近似）。
首次最差可见度发生在两套条纹互移半周期：
$\Delta\phi=\pi$，
即
\[
\boxed{\psi=\frac{\lambda}{2d}}
=\frac{0.55\times10^{-6}}{2(0.95)}
\simeq2.9\times10^{-7}\,\mathrm{rad}.
\]
用 $1$rad $=206265''$，
\[
\boxed{\psi\simeq0.060''}.
\]
\excheck 两套条纹来自\emph{不相干}恒星，
相互不能相干叠加振幅，但两套强度花样错位会降低总可见度。

\diffractionproblem{134}{210}
透射光栅周期
$d=1.50\,\mathrm{\mu m}$，
波长 $\lambda=530$nm。
求最高可见级主极大对应角散度
$D=\mathrm d\theta/\mathrm d\lambda$，
单位为角分/nm：
(a) 垂直入射；
(b) 入射角 $\theta_0=45^\circ$。
\exsolution
有向光栅方程
$d(\sin\theta-\sin\theta_0)=m\lambda$，
对波长微分：
\[
\boxed{|D|=\frac{|m|}{d|\cos\theta|}}.
\]
(a) 正入射最高级
$|m|=\lfloor d/\lambda\rfloor=2$，
$\sin\theta=2\lambda/d$，
所以
\[
\boxed{D_a=
\frac2{d\sqrt{1-(2\lambda/d)^2}}
\simeq6.5\,\text{角分/nm}}.
\]
(b) 因角度向另一侧延伸，最高绝对级是
$m=-4$，
$\sin\theta=\sin45^\circ-4\lambda/d$，
\[
\boxed{D_b=
\frac4{d\sqrt{1-(\sin45^\circ-4\lambda/d)^2}}
\simeq13\,\text{角分/nm}}.
\]
\excheck 将 rad/m 换成角分/nm：
乘 $3437.75\times10^{-9}$；不应仅报告无量纲13。

\diffractionproblem{135}{210}
波长为 $\lambda$ 的平面光垂直入射周期 $d$ 的光栅。
求主极大角散度 $D=\mathrm d\theta/\mathrm d\lambda$
用衍射角 $\theta$ 表示的公式。
\exsolution
第 $m$ 级光栅方程
$d\sin\theta=m\lambda$。
对 $\lambda$ 微分
$d\cos\theta\,\mathrm d\theta=m\,\mathrm d\lambda$，
从而
$D=m/(d\cos\theta)$。
再由 $m/d=\sin\theta/\lambda$，
得到
\[
\boxed{D=\frac{\tan\theta}{\lambda}}.
\]
\excheck 固定的是\emph{衍射级次} $m$；若改变级次而非波长，上式不是同一个偏导数。

\diffractionproblem{136}{210}
波长
$\lambda=589.0$nm 的光垂直入射光栅，
周期 $d=2.5\,\mathrm{\mu m}$，
总刻线数 $N=10000$。
求二级主极大的全角宽（两侧最近零点之间）。
\exsolution
二级方向满足
$\sin\theta_2=2\lambda/d$。
$N$ 缝干涉因子的最近零点在
$d\sin\theta=(2\pm1/N)\lambda$，
故峰两侧的全角宽
\[
\boxed{\Delta\theta\simeq
\frac{2\lambda}{Nd\cos\theta_2}
=\frac{2\lambda}
{Nd\sqrt{1-(2\lambda/d)^2}}}.
\]
代数值
$\Delta\theta\simeq5.3\times10^{-5}$rad，
\[
\boxed{\Delta\theta\simeq11''}.
\]
\excheck 题中“angular width”采用全主峰零点间宽度；若定义半宽则差一个因子2。

\diffractionproblem{137}{210}
证明光垂直入射宽度为 $\ell$ 的透射光栅时，
其最大分辨本领不能超过 $\ell/\lambda$。
\exsolution
波长分辨本领由瑞利判据给出
$\mathcal R=\lambda/\delta\lambda=mN$，
其中 $N$ 是照亮的刻线数、$m$ 为衍射级次。
若光栅周期为 $d$，则 $N\simeq\ell/d$。
另一方面衍射必须满足
$|m\lambda/d|\le1$，
即 $|m|\le d/\lambda$。
两者结合
\[
\boxed{\mathcal R_{\max}=mN
\le\frac d\lambda\frac\ell d
=\frac\ell\lambda}.
\]
\excheck 极限只是几何与衍射级次的上界，实际仪器还受光栅效率、缺陷、像差影响。

\diffractionproblem{138}{210}
用衍射光栅说明：
按瑞利判据刚能分辨的两条相邻谱线的频差
$\delta\nu$ 等于参与干涉的最远两振动传播时间差 $\Delta t$ 的倒数。
\exsolution
对第 $m$ 级，
在最远两光栅缝之间的路程差约
$\Delta\ell=N d\sin\theta
=N m\lambda$，
传播时延
$\Delta t=\Delta\ell/c
=Nm/\nu$。
瑞利判据的相邻谱线刚可分辨时，
$\delta\nu/\nu\simeq1/(mN)$，
故
\[
\boxed{\delta\nu=
\frac{\nu}{mN}
=\frac1{\Delta t}}.
\]
\excheck 对有限光栅应取 $N-1$ 个缝间隔，更精确的系数含 $N-1$；
书中公式为刻线多时 $N\gg1$ 的近似。

\diffractionproblem{139}{210}
两条相近的单色谱线波长分别为
$600.000\,\mathrm{nm}$ 与 $600.050\,\mathrm{nm}$，
垂直入射透射光栅。问在怎样的衍射角
$\theta$ 下恰好达到瑞利分辨判据？
\textbf{宽度版本差异：}多个可查的 Mir 等英文题面记光栅总宽
$\ell=10.0\,\mathrm{mm}$；
本项目此前记录的上传扫描底本题面记为
$\ell=10.0\,\mathrm m$。
两者相差1000倍，必须分别计算。
\exsolution
瑞利判据为 $\lambda/\delta\lambda=mN$。
另一方面正入射时 $d\sin\theta=m\lambda$，
刻线数 $N=\ell/d$。
将级次与刻线间距消去，得到
\[
 \boxed{\sin\theta
 =\frac{\lambda^2}{\ell\,\delta\lambda}},
 \qquad
 \lambda\simeq600.0\,\mathrm{nm},\quad
 \delta\lambda=0.050\,\mathrm{nm}.
\]
\textbf{标准英文版本 $\ell=10.0\,\mathrm{mm}$：}
\[
 \sin\theta=\frac{(600.0\times10^{-9})^2}
 {(10.0\times10^{-3})(0.050\times10^{-9})}
 =0.720,
 \qquad
 \boxed{\theta\simeq46.05^\circ\approx46^\circ}.
\]
\textbf{若严格按已记录扫描题面
$\ell=10.0\,\mathrm m$：}
\[
 \sin\theta=7.20\times10^{-4},\qquad
 \boxed{\theta\simeq0.0413^\circ}.
\]
\excheck
维数检查：$\lambda^2/(\ell\,\delta\lambda)$ 无量纲。
按标准英文版 $\ell=10.0\,\mathrm{mm}$ 可复现书末
$46^\circ$；按扫描底本记录 $10.0\,\mathrm m$ 不能。
公开对照可见
\texttt{studiestoday.com/ie-irodov-solutions-chapter-29-diffraction-of-light-part-2}。
这一修订只明确排定两种版本的主次，
不声称在本次重看了用户扫描页。

\diffractionproblem{140}{210}
正入射的透射光栅宽 $\ell=6.5\,\mathrm{cm}$，
每毫米刻有200条线。研究波长
$\lambda=670.8\,\mathrm{nm}$、
双线差 $\delta\lambda=0.015\,\mathrm{nm}$ 的谱线。
求：(a) 第几级起可分辨双线；
(b) 在波长约 $670\,\mathrm{nm}$ 处，
该光栅的最小理论可分辨波长差。
\textbf{单位版本差异：}本项目先前记录的扫描底本
第一处记成“$670.8\,\mathrm{mm}$”，
但其后文又明确要求在 $670\,\mathrm{nm}$ 附近讨论；
另有英文原题转录写作 $670.8\,\mathrm{nm}$。
\exsolution
刻线总数
\[
 N=\ell\frac{200}{\mathrm{mm}}
 =(65\,\mathrm{mm})(200/\mathrm{mm})
 =13000,
 \quad
 d=\frac{1}{200}\,\mathrm{mm}=5.0\,\mathrm{\mu m}.
\]
在第 $m$ 级满足 $\mathcal R=mN$，由
$mN\ge\lambda/\delta\lambda$ 得
\[
 m\ge\frac{670.8}{0.015(13000)}
 \simeq3.440,
 \quad\boxed{m_{\min}=4}.
\]
正入射真实传播级次满足
$m\lambda/d\le1$，
故 $m_{\max}=\lfloor d/\lambda\rfloor=7$。
最小可分辨谱线差
\[
 \boxed{\delta\lambda_{\min}
 =\frac{\lambda}{m_{\max}N}
 \simeq\frac{670\,\mathrm{nm}}{7(13000)}
 =0.00736\,\mathrm{nm}
 \simeq7.4\,\mathrm{pm}}.
\]
\excheck
若按“$670.8\,\mathrm{mm}$”字面处理，
单线波长大于 $5.0\,\mathrm{\mu m}$ 的光栅周期约
$1.34\times10^5$ 倍，不可能形成第一级光谱，
更不可能得到第四级；故记录为扫描版本单位冲突，
不能把毫米和纳米分支的条件混入同一推导。

\diffractionproblem{141}{210}
一片宽
$\ell=10$mm 的透射光栅，在垂直入射时，
钠黄双线 $589.0$nm 与 $589.6$nm 从第五级开始可被分辨。
估计：(a) 光栅周期；
(b) 若仍用同样周期，欲在第三级分辨波长
$\lambda=460.0$nm、间隔
$\delta\lambda=0.13$nm 的双线，光栅宽度至少应取多少？
\exsolution
(a) 临界分辨本领
$mN\simeq\lambda/\delta\lambda$，
取 $m=5$，
$N\simeq589/[5(0.6)]\simeq196$。
周期
\[
\boxed{d\simeq\frac{\ell}{N}
\simeq\frac{10\,\mathrm{mm}}{196}
\simeq0.050\,\mathrm{mm}}.
\]
(b) 第三级
$N_{\rm req}\simeq460/[3(0.13)]
\simeq1180$，
所以
\[
\boxed{\ell_{\min}\simeq N_{\rm req}d
\simeq1180(0.050\,\mathrm{mm})
\simeq6.0\,\mathrm{cm}}.
\]
\excheck 第一级别能分辨是在$m=5$，
并不说明其他更高衍射级的主峰强度足够强；
此处只比较理论分辨本领。

\diffractionproblem{142}{210--211}
石英光谱仪透射光栅宽
$\ell=25$mm、每毫米250条线，
物镜焦距
$f=80$cm。
波长为 $310.154$ 与
$310.184$nm 的双线垂直入射。
求：(a) 在一级、二级谱中像面上的分离距离；
(b) 两级是否可分辨。
\exsolution
光栅周期
$d=1/250\,\mathrm{mm}=4\,\mathrm{\mu m}$，
总刻线数
$N=25(250)=6250$。
双线差
$\delta\lambda=0.030$nm，
第 $m$ 级满足
$d\sin\theta_m=m\lambda$。
焦平面坐标 $x_m=f\tan\theta_m$，
故
\[
\delta x_m\simeq
f\frac{m\,\delta\lambda}{d\cos^3\theta_m}.
\]
将 $\lambda_1=310.154\,\mathrm{nm}$、
$\lambda_2=310.184\,\mathrm{nm}$ 分别代入
$x_m(\lambda)=f\tan[\arcsin(m\lambda/d)]$，
可避免有限谱线间隔的线性化误差，得到
\[
 \boxed{\delta x_1\simeq6.055\,\mathrm{\mu m}},
 \qquad
 \boxed{\delta x_2\simeq12.446\,\mathrm{\mu m}}.
\]
按近轴近似 $\cos\theta_m\simeq1$，
有 $\delta x_m\simeq fm\,\delta\lambda/d$，
对应常用有效数字结果
$\boxed{6\,\mathrm{\mu m},\quad12\,\mathrm{\mu m}}$；
二级相对精确焦平面值偏低约 $3.7\%$。
瑞利分辨条件
$\delta\lambda\ge\lambda/(mN)$：
一级阈值约 $0.050$nm（大于0.030nm），不可分辨；
二级约 $0.025$nm（小于0.030nm），可以分辨。
\[
\boxed{\text{一级不能分辨；二级可以分辨。}}
\]
\excheck 判断能否分辨应使用相同的光栅线数与瑞利准则，而不是只看像面双线分离几微米。
% === END 29-optics-diffraction-131-142.tex ===

% === BEGIN 29-optics-diffraction-143-150.tex blob 3b54d3d1b3ae21cdc152273373af4b10beb0cb78 ===
% Irodov Part V §5.3 originals 143--150; printed pp.211--212.
% Original book answers pp.331--332.
\diffractionproblem{143}{211}
光谱仪的三棱镜在最小偏折状态，底边宽度为 $b$（图5.28），
材料折射率 $n=n(\lambda)$。
利用夫琅禾费衍射与瑞利分辨准则证明极限分辨本领
$\lambda/\delta\lambda=b|\mathrm dn/\mathrm d\lambda|$。
\exsolution
考虑来自棱镜左右极端的两条光线，
它们的光程差在最小偏折方向满足同一波前条件。
波长改变 $\delta\lambda$ 引起折射率改变量
$\delta n=(\mathrm dn/\mathrm d\lambda)\delta\lambda$，
于是两极端出射波新增相位对应光程差
$\delta\Delta=b\,\delta n$。
根据瑞利准则，当一条谱线的中心极大与另一谱线的第一个衍射极小重合时，
两极端新增的光程差约等于一个波长：
$|b\,\delta n|=\lambda$。
因此
\[
b\left|\frac{\mathrm dn}{\mathrm d\lambda}\right|
\delta\lambda\simeq\lambda,
\qquad
\boxed{\frac{\lambda}{\delta\lambda}
=b\left|\frac{\mathrm dn}{\mathrm d\lambda}\right|}.
\]
\excheck 这是假设棱镜的分辨受其底边有效光束口径限制的\emph{衍射极限}；与仅给角散度的公式不同。

\diffractionproblem{144}{211}
光谱仪的三棱镜玻璃色散满足
$n=A+B/\lambda^2$，其中
$B=0.010\,\mathrm{\mu m^2}$。
(a) 求极限分辨本领随波长的关系；棱镜底边
$b=5.0$cm 时在
$\lambda_1=434$nm、
$\lambda_2=656$nm 处的数值；
(b) 分辨钠黄双线
$589.0,589.6$nm 所需的底边宽度。
\exsolution
对色散关系求导
$\mathrm dn/\mathrm d\lambda=-2B/\lambda^3$。
由第143题
\[
\boxed{\mathcal R=\frac{\lambda}{\delta\lambda}
=\frac{2bB}{\lambda^3}},
\]
说明 $\mathcal R\propto\lambda^{-3}$。
单位一致地取 $b=50000\,\mathrm{\mu m}$：
\[
\boxed{\mathcal R(0.434\,\mathrm{\mu m})
\simeq1.2\times10^4},\quad
\boxed{\mathcal R(0.656\,\mathrm{\mu m})
\simeq0.35\times10^4}.
\]
(b) 分辨条件
$2bB/\lambda^3\ge\lambda/\delta\lambda$，
所以
\[
\boxed{b_{\min}=
\frac{\lambda^4}{2B\delta\lambda}}.
\]
取 $\lambda\simeq0.589\,\mathrm{\mu m}$、
$\delta\lambda=0.0006\,\mathrm{\mu m}$，
$\boxed{b_{\min}\simeq1.0\,\mathrm{cm}}$。
\excheck $B$ 的单位为 $\mathrm{\mu m^2}$；
若把 $b$ 的厘米数直接当微米使用，会使分辨本领差四个数量级。

\diffractionproblem{145}{211}
若棱镜在目标波长附近
$|\mathrm dn/\mathrm d\lambda|
=0.10\,\mathrm{\mu m^{-1}}$，
求棱镜底边宽度，
使它具有一片刻线数
$N=10000$ 的光栅在二级谱中的相同分辨本领。
\exsolution
二级光栅分辨本领
$\mathcal R_{\rm gr}=mN=2(10000)=20000$。
棱镜衍射极限分辨本领
$\mathcal R_{\rm prism}=b|\mathrm dn/\mathrm d\lambda|$。
令二者相同：
\[
\boxed{b=
\frac{mN}{|\mathrm dn/\mathrm d\lambda|}
=\frac{20000}{0.10}\,\mathrm{\mu m}
=0.20\,\mathrm m=20\,\mathrm{cm}}.
\]
\excheck 这里采用棱镜\emph{底边}长度而非光路宽的任意投影，光栅亦以实际被照亮的刻线数计算。

\diffractionproblem{146}{211}
物镜直径 $D=5.0$cm 的望远镜，取
$\lambda=0.55\,\mathrm{\mu m}$。
求物镜的角分辨本领以及距望远镜
$\ell=3.0$km 的两物点的最小可分辨间距。
\exsolution
圆孔艾里斑瑞利判据
$\delta\psi_{\min}=1.22\lambda/D$，
角分辨本领
\[
\boxed{\mathcal R_\psi=
\frac1{\delta\psi_{\min}}
=\frac D{1.22\lambda}
\simeq7.4\times10^4}.
\]
远处两点最小间距
$\Delta y\simeq\ell\delta\psi_{\min}$，
所以
\[
\boxed{\Delta y_{\min}
\simeq\ell\frac{1.22\lambda}{D}
=\frac{3000(1.22)(0.55\times10^{-6})}{0.050}
\simeq4.0\,\mathrm{cm}}.
\]
\excheck 这里的“分辨本领”按\emph{最小角度的倒数}定义，而不是线性放大率。

\diffractionproblem{147}{211}
口径 $D=5$m 的反射望远镜在
$\lambda=0.55\,\mathrm{\mu m}$ 处工作。
求在月面上能分辨的两个点之间的最小距离。
\exsolution
由瑞利判据最小角距
$\delta\psi=1.22\lambda/D$。
取平均地月距离
$L_{\rm Moon}\simeq3.84\times10^8$m，
则
\[
\boxed{\Delta y_{\min}=
L_{\rm Moon}\delta\psi
=\frac{(3.84\times10^8)(1.22)(0.55\times10^{-6})}{5}
\simeq5.2\times10^1\,\mathrm m
\approx50\,\mathrm m}.
\]
\excheck 此为理想衍射极限；地面大气视宁度通常会劣化实际角分辨率，不能直接当成现实观测保证。

\diffractionproblem{148}{211}
望远镜物镜口径
$D=5.0$cm，观察者瞳孔直径
$d_0=4.0$mm。
求完全利用物镜理论分辨本领所需的最小角放大率。
\exsolution
物镜的瑞利最小角距
$\delta\psi=1.22\lambda/D$，
人眼瞳孔对同波长的理想衍射分辨极限
$\delta\psi_{\rm eye}=1.22\lambda/d_0$。
望远镜使角距放大 $\Gamma$，
使
$\Gamma\delta\psi\ge\delta\psi_{\rm eye}$。
因此
\[
\boxed{\Gamma_{\min}=
\frac{\delta\psi_{\rm eye}}{\delta\psi}
=\frac D{d_0}
=\frac{50\,\mathrm{mm}}{4.0\,\mathrm{mm}}
=12.5\approx13}.
\]
\excheck 这里的最小放大率是衍射分辨利用条件，
不是在所有目标亮度下的最佳视觉放大率。

\diffractionproblem{149}{212}
显微镜物镜的数值孔径
$\mathrm{NA}=\sin\alpha=0.24$，
用波长
$\lambda=0.55\,\mathrm{\mu m}$ 的光照亮物体。
求该显微镜能分辨的两物点最小间距。
\exsolution
圆孔物镜在物方的瑞利横向分辨极限
\[
\boxed{d_{\min}=
\frac{0.61\lambda}{\mathrm{NA}}
=\frac{0.61(0.55\,\mathrm{\mu m})}{0.24}
\simeq1.4\,\mathrm{\mu m}}.
\]
\excheck 原题给的数值孔径已是 $\sin\alpha$，
不可再对0.24取一次正弦；
这里默认物体位于空气中，折射率为1。

\diffractionproblem{150}{212}
显微镜物镜数值孔径为
$\sin\alpha=0.24$，观察者瞳孔
$d_0=4.0$mm。
求使物镜衍射分辨本领被眼睛完全利用的最小放大率。
\exsolution
显微镜物方线分辨极限
$d_{\min}=0.61\lambda/\sin\alpha$。
若将它放在标准明视距离
$\ell_0=25$cm 观看，
该间距所张裸眼角为
$d_{\min}/\ell_0$。
眼瞳衍射角极限约
$1.22\lambda/d_0$。
需放大
$\Gamma_{\min}d_{\min}/\ell_0
=1.22\lambda/d_0$，故
\[
\boxed{\Gamma_{\min}=
\frac{2\ell_0\sin\alpha}{d_0}
=\frac{2(0.25)(0.24)}{0.004}
=30}.
\]
\excheck 公式中波长消掉，这只是将物镜极限角尺度匹配到眼瞳的衍射极限；视网膜和视觉神经等实际因素未纳入。
% === END 29-optics-diffraction-143-150.tex ===

% === BEGIN 29-optics-diffraction-151-156.tex blob b624f1c4e27f36f4c58e811e9b24209dbea43755 ===
% Irodov Part V §5.3 originals 151--156 (end of diffraction section).
% Uploaded PDF printed pp.212; original answers pp.331--332.
\diffractionproblem{151}{212}
X 射线以掠入射角
$\alpha_0=60.0^\circ$ 照射周期为 $a$ 的一维线性散射中心链。
若 $\lambda=2a/5$，
求所有出现衍射主极大的入射或散射方向角。
\exsolution
以原书把光与链轴夹角记为 $\alpha$，
相邻散射中心的相位差（以入射方向60°为基准）为
$(2\pi a/\lambda)(\cos\alpha-\cos60^\circ)$。
相长需它为 $2\pi m$：
\[
\boxed{\cos\alpha=\cos60^\circ+m\frac{\lambda}{a}
=\frac12+\frac{2m}{5}}.
\]
要求 $-1\le\cos\alpha\le1$，
给整数
$m=-3,-2,-1,0,1$。
逐一计算：
\[
\boxed{\alpha\simeq
134^\circ,\ 107^\circ,\ 84^\circ,\ 60^\circ,\ 26^\circ}.
\]
\excheck 题目提到的是与\emph{链轴的掠入射方向角}，不是普通二维平面光栅的法线入射角；必须用余弦方向差。

\diffractionproblem{152}{212}
波长
$\lambda=40$pm 的 X 射线垂直照到二维矩形散射点阵（图5.29）。
屏距 $\ell=10$cm，
$x$ 轴两个对称二级极大间隔
$\Delta x=60$mm，
$y$ 轴两个对称二级极大间隔
$\Delta y=40$mm。
求两个点阵周期 $a,b$。
\exsolution
两个对称二级极大各离中心的半距为
$x=\Delta x/2$、
$y=\Delta y/2$。
横向出射角满足
$\tan\theta_x=x/\ell$，
$\tan\theta_y=y/\ell$。
二级点阵相长条件
$a\sin\theta_x=2\lambda$、
$b\sin\theta_y=2\lambda$。
故
\[
\boxed{a=\frac{2\lambda}
{\sin(\arctan[\Delta x/(2\ell)])}
\simeq0.28\,\mathrm{nm}},
\]
\[
\boxed{b=\frac{2\lambda}
{\sin(\arctan[\Delta y/(2\ell)])}
\simeq0.41\,\mathrm{nm}}.
\]
\excheck 图5.29给的是对称极大之间的\emph{总间距}，计算散射角时必须先除以2。

\diffractionproblem{153}{212}
一束 X 射线沿三维直角点阵的 $a$ 轴周期方向入射；
另外两条互相垂直的晶格周期为 $b,c$。
求衍射主极大方向和可能的波长。
\exsolution
设出射方向与三轴的夹角为
$\alpha,\beta,\gamma$，
方向余弦满足
$\cos^2\alpha+\cos^2\beta+\cos^2\gamma=1$。
对三方向相邻散射中心的相位差应用相长条件，得
\[
\boxed{a(1-\cos\alpha)=m_1\lambda},\quad
\boxed{b\cos\beta=m_2\lambda},\quad
\boxed{c\cos\gamma=m_3\lambda},
\]
其中 $m_i$ 为使所给方向物理可行的整数。
故
$\cos\alpha=1-m_1\lambda/a$、
$\cos\beta=m_2\lambda/b$、
$\cos\gamma=m_3\lambda/c$。
代入单位方向约束，消去二次项中常数：
\[
-2\frac{m_1\lambda}{a}
+\lambda^2\left[
\left(\frac{m_1}{a}\right)^2+
\left(\frac{m_2}{b}\right)^2+
\left(\frac{m_3}{c}\right)^2
\right]=0.
\]
非零波长的解
\[
\boxed{\lambda=
\frac{2m_1/a}
{(m_1/a)^2+(m_2/b)^2+(m_3/c)^2}}.
\]
\excheck 选取整数三元组后仍必须核查
$|\cos\alpha|,|\cos\beta|,|\cos\gamma|\le1$ 与 $\lambda>0$。

\diffractionproblem{154}{212}
X 射线以掠射角
$\alpha=60.0^\circ$ 照射氯化钠单晶天然晶面，
NaCl 密度
$\rho=2.16\,\mathrm{g/cm^3}$，
镜面反射为二级布拉格极大。
求辐射波长。
\exsolution
设一个 NaCl 化学式单位的质量
$m=M/N_A$，
$M=58.44\,\mathrm{g/mol}$。
NaCl 岩盐晶胞包含4个化学式单位，
晶胞边长
$a_{\rm lat}=(4m/\rho)^{1/3}$。
该天然晶面相邻反射层间距
$d=a_{\rm lat}/2=(m/(2\rho))^{1/3}$。
布拉格方程
$2d\sin\alpha=k\lambda$，
$k=2$，所以
\[
\boxed{\lambda=\frac{2}{k}
\left(\frac{m}{2\rho}\right)^{1/3}
\sin\alpha}.
\]
代
$\rho=2160\,\mathrm{kg/m^3}$、
$m=0.05844/(6.02214076\times10^{23})$kg，
\[
\boxed{\lambda\simeq244\,\mathrm{pm}}.
\]
\excheck 应使用NaCl岩盐结构的反射层间距，
而非把晶胞边长直接当布拉格层距；题目没有单独给晶格常数，须由密度和单胞粒子数求出。

\diffractionproblem{155}{212}
波长
$\lambda=174$pm 的 X 射线入射一旋转单晶，
晶体围绕平行表面且垂直入射方向的轴旋转。
平行表面的晶面组的二级与三级衍射峰方向夹角
$\alpha=60^\circ$。
求晶面间距 $d$。
\exsolution
布拉格条件
$2d\sin\theta_2=2\lambda$、
$2d\sin\theta_3=3\lambda$。
由于不同衍射级的反射方向在旋转扫描时夹角为
$\alpha$，对应掠射角的差为
$|\theta_3-\theta_2|=\alpha/2$。
利用
$\sin(\theta_3-\theta_2)
=\sin\theta_3\cos\theta_2-\cos\theta_3\sin\theta_2$，
结合两式消元，可得
\[
\boxed{d=
\frac{\lambda}{2\sin(\alpha/2)}
\sqrt{k_2^2+k_3^2-2k_2k_3\cos(\alpha/2)}},
\quad k_2=2,\ k_3=3.
\]
取 $\alpha/2=30^\circ$：
\[
\boxed{d\simeq
\frac{174\,\mathrm{pm}}{1}
\sqrt{13-12\cos30^\circ}
\simeq281\,\mathrm{pm}
=0.28\,\mathrm{nm}}.
\]
\excheck 书末答案的纯文本提取似为“0.28 pm”，但公式与174pm入射波长严格给\emph{0.28 nm}；单位需在勘误注明。

\diffractionproblem{156}{212}
波长
$\lambda=17.8$pm 的 X 射线穿过多晶试样，
距样品 $\ell=15$cm 的屏上出现环形衍射条纹。
晶面间距
$d=155$pm，
求第二级布拉格反射对应的亮环半径。
\exsolution
第二级布拉格条件
$2d\sin\alpha=k\lambda$，
$k=2$，所以
$\sin\alpha=2\lambda/(2d)=\lambda/d$，
$\alpha\simeq6.59^\circ$。
衍射光束与直射光束的夹角为
$2\alpha$，
环半径满足
$r=\ell\tan(2\alpha)$。
故
\[
\boxed{r=\ell
\tan\left[2\arcsin\frac{k\lambda}{2d}\right]
\simeq0.15\tan(13.18^\circ)
\simeq3.5\,\mathrm{cm}}.
\]
\excheck 多晶体德拜环几何出射角为 $2\alpha$，
而不是直接把布拉格掠射角 $\alpha$ 代入屏上半径公式。
% === END 29-optics-diffraction-151-156.tex ===
